% ===========================================================================
%  Quality Inspector Portal - standalone manual
%  Build:  pdflatex manual-qc.tex      (run it twice so the contents page is right)
% ===========================================================================
\documentclass[11pt,a4paper,oneside]{report}
\usepackage{style/manualstyle}

\hypersetup{pdftitle={Denim RFID Track and Trace - Quality Inspector Portal}}

\begin{document}

\begin{titlepage}
\begin{tikzpicture}[remember picture,overlay]
  \fill[secQC] (current page.north west) rectangle ([yshift=-95mm]current page.north east);
  \node[anchor=north west,text=white,font=\fontsize{30}{36}\bfseries,align=left]
    at ([xshift=25mm,yshift=-26mm]current page.north west)
    {Denim RFID\\Track \& Trace};
  \node[anchor=north west,text=white,font=\LARGE,align=left]
    at ([xshift=25mm,yshift=-66mm]current page.north west) {Quality Inspector Portal};
  \node[anchor=south east,text=brandgrey,font=\small,align=right]
    at ([xshift=-25mm,yshift=25mm]current page.south east) {Inspecting garments \textbullet\ Marking faults on the design\\Passing and failing \textbullet\ Defect analysis};
\end{tikzpicture}

\vspace*{110mm}

{\Large\bfseries For staff at the inspection benches}

\vspace{6mm}

{\large\color{brandgrey}
This booklet assumes no computer experience. Every screen you will meet is
pictured, and every button you need to press is named.}

\vspace{14mm}

\begin{plainbox}
\textbf{What this software does, in one sentence.}\\[4pt]
Every pair of jeans gets a small electronic tag when it is stitched, and from
that moment the computer always knows where that garment is, how long it has
been there, and who last handled it.
\end{plainbox}

\end{titlepage}

\tableofcontents
\clearpage

\chapter{About This Software}

\section{What problem does it solve?}

A large denim factory makes tens of thousands of garments every day. Those
garments are constantly moving: from the cutting room to the sewing lines, into
the washing plant, on to finishing, through quality checking, and finally into
cartons for the customer.

The old way of keeping track was paper. Somebody counted a bundle, wrote the
number on a slip, and passed it along. When the numbers did not match at the end
of the day, nobody could say where the missing pieces went.

This software replaces that guesswork. A small electronic tag --- called an
\textbf{RFID tag} --- is attached to every single garment. Special readers can
sense hundreds of these tags at once, without anyone having to look at them or
scan them one by one.

\begin{plainwords}
\textbf{RFID} stands for ``Radio Frequency Identification''. Think of it as a
tiny electronic name badge. It has no battery. When a reader sends out a radio
signal, the tag wakes up and calls back its own unique number. A reader can hear
hundreds of tags in a couple of seconds, even inside a closed trolley.
\end{plainwords}

\section{What the system keeps track of}

\begin{bullets}
  \item \textbf{Fabric rolls} arriving from suppliers, and how much of each roll has been used.
  \item \textbf{Bundles} of cut pieces going from the cutting room to the sewing lines.
  \item \textbf{Every finished garment}, individually, from stitching to shipping.
  \item \textbf{Every handover} between departments, with a document and a count check.
  \item \textbf{Every quality inspection}, including exactly where on the garment a fault was found.
  \item \textbf{Every person} who touched the work, and when.
\end{bullets}

\section{The journey of one pair of jeans}

Figure~\ref{fig:flow} shows the whole route. Each box is a department. Each
arrow is a handover, and at every handover the system does the same two things:
the department sending the goods scans what is leaving, and the department
receiving them scans what arrived. The computer then compares the two.

\begin{figure}[H]
\centering
\begin{tikzpicture}[
  node distance=6mm and 9mm,
  box/.style={rectangle,rounded corners=3pt,draw=none,text=white,
              minimum height=11mm,minimum width=27mm,align=center,font=\small\bfseries},
  arrow/.style={-{Stealth[length=3mm]},thick,draw=brandgrey},
  lbl/.style={font=\scriptsize\itshape,text=brandgrey}]

\node[box,fill=secFabric] (fab) {Fabric\\Warehouse};
\node[box,fill=secCut,right=of fab] (cut) {Cutting};
\node[box,fill=secStitch,right=of cut] (sti) {Stitching};
\node[box,fill=secSort,right=of sti] (sor) {Sorting};

\node[box,fill=secWash,below=14mm of sor] (was) {Washing \&\\Treatment};
\node[box,fill=secFinish,left=of was] (fin) {Finishing};
\node[box,fill=secQC,left=of fin] (qc) {Quality\\Control};
\node[box,fill=secRetro,left=of qc] (ret) {Retrofitting};

\node[box,fill=secDisp,below=14mm of ret] (dis) {Dispatch \&\\Packing};
\node[box,fill=brandgrey,right=of dis] (shp) {Shipped};

\draw[arrow] (fab) -- node[lbl,above]{rolls} (cut);
\draw[arrow] (cut) -- node[lbl,above,align=center]{hand count} (sti);
\draw[arrow] (sti) -- node[lbl,above,align=center]{tag added} (sor);
\draw[arrow] (sor) -- (was);
\draw[arrow] (was) -- (fin);
\draw[arrow] (fin) -- (qc);
\draw[arrow] (qc) -- node[lbl,above]{failed} (ret);
\draw[arrow] (ret) -- (dis);
\draw[arrow] (dis) -- node[lbl,above,align=center]{new tag} (shp);

\draw[arrow,dashed] (ret.north) to[out=60,in=200] node[lbl,above,sloped]{corrected, back to QC} (qc.south west);
\draw[arrow,dashed] (qc.south) to[out=280,in=20] node[lbl,below,sloped]{passed} (dis.north east);

\end{tikzpicture}
\caption{The route every garment takes. Solid arrows are the normal path;
dashed arrows show what happens after a quality check.}
\label{fig:flow}
\end{figure}

\section{The one idea you must understand: the two-sided count}

This is the heart of the system, and it is simple.

\begin{figure}[H]
\centering
\begin{tikzpicture}[
  node distance=10mm,
  dept/.style={rectangle,rounded corners=3pt,fill=brandsoft,draw=brandblue,
               minimum height=16mm,minimum width=38mm,align=center,font=\small\bfseries},
  doc/.style={rectangle,rounded corners=2pt,fill=white,draw=brandgrey,
              minimum height=13mm,minimum width=34mm,align=center,font=\scriptsize},
  arrow/.style={-{Stealth[length=3mm]},thick,draw=brandblue}]

\node[dept] (a) {Sending\\department};
\node[doc,right=22mm of a] (d) {Transfer note\\\textbf{240 pieces}};
\node[dept,right=22mm of d] (b) {Receiving\\department};

\draw[arrow] (a) -- node[font=\scriptsize,above,align=center]{scans what\\is leaving} (d);
\draw[arrow] (d) -- node[font=\scriptsize,above,align=center]{scans what\\arrived} (b);

\node[below=9mm of d,font=\small,align=center,text=brandink]
  {The computer compares the two numbers};
\node[below=17mm of d,align=center,font=\small\bfseries]
  {\textcolor{okgreen}{Same $\Rightarrow$ MATCHED}\qquad
   \textcolor{dangerred}{Different $\Rightarrow$ VARIANCE}};
\end{tikzpicture}
\caption{Every handover between departments works this way.}
\label{fig:tally}
\end{figure}

If the two counts agree, the handover is marked \textbf{MATCHED} and everyone
moves on. If they do not agree, the system marks it a \textbf{VARIANCE} and
shows a list of exactly which garments were not found --- by serial number, style
and size. Nobody has to hunt through paperwork.

\begin{tipbox}
A variance is not an accusation. Tags are sometimes missed because a garment was
at the bottom of a trolley. Simply scan the batch again; anything found the
second time is added, and the variance clears by itself.
\end{tipbox}

\section{Who uses which portal}

The same software shows a different set of screens depending on who signs in.
This manual has one chapter for each group.

\begin{table}[H]
\centering
\small
\begin{tabular}{@{}p{34mm}p{45mm}p{60mm}@{}}
\toprule
\textbf{Portal} & \textbf{Who signs in} & \textbf{What they mainly do} \\
\midrule
Operator &
  Machine and station staff in the store, cutting, sewing, sorting, wash,
  finishing, retrofit and dispatch &
  Scan garments in and out, attach tags, count bundles, correct faults \\
Quality Inspector &
  QC staff at the inspection benches &
  Check garments, mark faults on the picture of the design, pass or fail \\
Supervisor &
  The person in charge of a department &
  Everything an operator can do, plus settle count differences and cancel
  mistaken handovers \\
Administrator &
  Plant management and the IT administrator &
  See all the numbers, build reports, manage staff accounts, register readers,
  maintain the product list \\
\bottomrule
\end{tabular}
\caption{The four portals covered by this manual.}
\end{table}

\chapter{Getting Started}

\section{What you need}

\begin{bullets}
  \item A computer or tablet with a web browser (Chrome, Edge or Firefox).
  \item The address of the system. On the factory network this is usually
        \typed{http://denim-rfid:8080} or an address your supervisor will give you.
  \item Your own username and password. Never use somebody else's.
\end{bullets}

\begin{stopbox}
Your username identifies every action you take. If you scan a batch under
somebody else's login, the record will say they did it. Always sign in as
yourself, and sign out when you leave the station.
\end{stopbox}

\section{Signing in}

\begin{steps}
  \item Open the browser and type the address of the system.
  \item Type your \field{Username} in the first box.
  \item Type your \field{Password} in the second box.
  \item Press \btn{Sign in}.
\end{steps}

\shotsmall{01-login}{The sign-in screen. This is the first thing you see.}{fig:login}

If the password is wrong you will see a red message. Check that Caps~Lock is
off and try again. After several failed attempts, ask your supervisor --- they
can set you a new password.

\section{Signing out}

Press \btn{Sign out} at the bottom left of the screen. Always do this before
leaving a shared station, otherwise the next person's work is recorded under
your name.

\section{The parts of the screen}

Every screen in the system is laid out the same way. Once you know these four
areas, you can find your way around any part of the software.

\shot{11-admin-dashboard-top}{The four areas of every screen: (1) the menu down
the left, (2) your name and the sign-out button at the bottom left, (3) the page
title and buttons across the top, (4) the working area in the middle.}{fig:layout}

\begin{table}[H]
\centering
\small
\begin{tabular}{@{}p{30mm}p{110mm}@{}}
\toprule
\textbf{Area} & \textbf{What it is for} \\
\midrule
Left menu &
  The list of screens you are allowed to open. You will only see the ones your
  job needs --- the list is shorter for an operator than for a manager. A number
  in a circle means something is waiting for you. \\
Bottom left &
  Your name, your role, the \btn{Theme} button (switches between light and dark
  for bright or dim areas) and \btn{Sign out}. \\
Top bar &
  The name of the screen you are on, a one-line description, and the main
  buttons for that screen. \\
Middle &
  The actual work: lists, forms and the boxes you scan into. \\
\bottomrule
\end{tabular}
\caption{What each part of the screen does.}
\end{table}

\begin{tipbox}
If the left menu is missing on a small tablet, press the \btn{$\equiv$} button at
the top left to slide it open.
\end{tipbox}

\chapter{Things Every User Needs to Know}

This chapter explains the handful of ideas that appear on nearly every screen.
Read it once and the rest of the manual will make sense.

\section{The scan box}

Wherever the system needs to know which garments you are holding, you will see a
large box like the one in Figure~\ref{fig:scanbox}. This is the \textbf{scan box}.

\shotsmall{71-transfers-receive}{The scan box. The big number on the blue strip
counts the tags you have read so far.}{fig:scanbox}

There are three ways tags get into that box.

\begin{steps}
  \item \textbf{A handheld reader.} Most hand scanners are set up to behave like
        a keyboard. Click inside the box once, then pull the trigger. The tag
        numbers type themselves in, one per line.
  \item \textbf{A fixed reader.} At a portal or tunnel the reader sends the tags
        to the system by itself. You do not have to do anything.
  \item \textbf{The \btn{Simulate reader} button.} This is for training and
        testing. It fills the box with the tags that really are in that location,
        so you can practise without hardware.
\end{steps}

\begin{tipbox}
Reading the same tag twice does no harm. The system counts each garment once,
however many times its tag is seen. If in doubt, scan again.
\end{tipbox}

The blue strip shows how many tags are ready. \btn{Clear} empties the box and
starts again.

\section{Coloured labels}

Small coloured labels appear all over the system. The colour tells you the
meaning at a glance.

\begin{table}[H]
\centering
\small
\begin{tabular}{@{}p{26mm}p{26mm}p{86mm}@{}}
\toprule
\textbf{Colour} & \textbf{Examples} & \textbf{Meaning} \\
\midrule
Green   & \texttt{PASS}, \texttt{RECEIVED}, \texttt{READY}
        & Everything is fine. Nothing to do. \\
Blue    & \texttt{IN TRANSIT}, \texttt{DISPATCHED}
        & Normal, in progress. Somebody else has the next step. \\
Amber   & \texttt{VARIANCE}, \texttt{REWORK}, \texttt{HOLD}
        & Needs a person to look at it. \\
Red     & \texttt{FAIL}, \texttt{SCRAP}, \texttt{MISSING}
        & Something is wrong, or the garment has been rejected. \\
Grey    & \texttt{PENDING}, \texttt{CUT}
        & Waiting its turn. Normal. \\
\bottomrule
\end{tabular}
\caption{What the colours mean.}
\end{table}

\section{Words you will see}

\begin{table}[H]
\centering
\small
\begin{tabular}{@{}p{34mm}p{104mm}@{}}
\toprule
\textbf{Word} & \textbf{What it means in plain language} \\
\midrule
Tag / EPC &
  The electronic label on the garment, and the unique number stored inside it.
  \texttt{EPC} is just the technical name for that number. \\
Article &
  One single garment. Every article has its own tag and its own serial number. \\
Serial number &
  A readable name for a garment, like \typed{ART-260816-000123}. Easier to say
  aloud than the tag number. \\
Bundle &
  A tied stack of cut pieces --- for example fifty size 32 fronts and backs ---
  handed from cutting to the sewing line. \\
Transfer note &
  The document created when one department sends goods to another. It lists
  exactly what was sent. \\
Batch &
  A group of garments moved together, usually all the same design and size. \\
Variance &
  The count on arrival did not match the count that was sent. \\
WIP &
  ``Work in progress'': everything currently on the factory floor, not yet shipped. \\
Commissioning &
  The moment a tag is attached to a garment and registered in the system.
  After this, the garment exists as far as the computer is concerned. \\
\bottomrule
\end{tabular}
\caption{The vocabulary of the system.}
\end{table}

\section{Tables: sorting, searching and opening}

Most screens show a table. Three things are always true:

\begin{bullets}
  \item \textbf{Rows can be clicked.} Clicking a row opens the full detail of that item.
  \item \textbf{Counts are on the right.} Numbers line up so you can compare them down a column.
  \item \textbf{\btn{Export CSV} saves the table} to a file you can open in Excel.
\end{bullets}

\section{If something goes wrong}

The system tells you when it cannot do what you asked, in a small panel at the
bottom right of the screen. It always explains why, in words --- for example
``Currently in Sorting Station, not Washing''.

\begin{warnbox}
The system refuses things on purpose to protect the count. If a garment is
rejected, it is nearly always because it is genuinely somewhere else, or its tag
belongs to another garment. Put that piece aside and tell your supervisor rather
than trying again repeatedly.
\end{warnbox}

Chapter~\ref{ch:trouble} lists the messages you are most likely to meet and what
to do about each one.

\section{Moving goods between departments}

Every station does these two things, whatever else its job is. Read this once;
the rest of the manual refers back to it.

\section{Receiving goods into your section}
\label{sec:receive}

Open \menu{Transfers} from the left menu. The top table, \emph{Waiting to be
received}, lists everything other departments have sent you.

\shot{70-transfers-inbox}{The Transfers screen. Everything sent to your section
is at the top, with how long it has been waiting.}{fig:op-inbox}

The \field{Waiting} column turns amber after four hours and red after eight.
That is your cue to deal with it.

\begin{steps}
  \item Find the row for the trolley in front of you. The \field{Document}
        number is printed on the note that came with it.
  \item Press \btn{Receive} on that row.
  \item Bulk-read everything that arrived into the scan box.
  \item Press \btn{Receive and tally}.
\end{steps}

\subsection{Reading the result}

The system now compares what arrived against what was sent, and shows five boxes.

\shotsmall{72-transfers-tally}{The result of a count. Here two pieces were not
found, so the handover is marked as a variance and the missing garments are
listed by name.}{fig:op-tally}

\begin{table}[H]
\centering
\small
\begin{tabular}{@{}p{24mm}p{110mm}@{}}
\toprule
\textbf{Box} & \textbf{Meaning} \\
\midrule
Expected & How many the sending department said they sent. \\
Received & How many you have just found. \\
Missing  & On the note but not found. \\
Extra    & Found but not on the note. \\
Result   & \texttt{MATCHED} if everything agrees, \texttt{VARIANCE} if not. \\
\bottomrule
\end{tabular}
\caption{The five boxes shown after a count.}
\end{table}

\begin{tipbox}
\textbf{If it says VARIANCE, do not panic.} Look through the trolley again and
scan a second time. Anything found is added to the count and the variance clears
by itself. Only if the pieces genuinely cannot be found does a supervisor need to
close the variance --- your supervisor has to close it.
\end{tipbox}

\section{Sending goods to the next department}
\label{sec:dispatch}

\begin{steps}
  \item On the \menu{Transfers} screen press \btn{+ Dispatch a batch}.
  \item Choose the department in \field{Send to}.
  \item Type a \field{Batch reference} if your section uses one, for example
        \typed{WASH-LOT-07}. This is printed on the note so the receiving
        department can find it.
  \item Bulk-read everything that is leaving.
  \item Press \btn{Check what was read} if you want to see a breakdown first.
  \item Press \btn{Generate transfer note}.
\end{steps}

\shotsmall{73-transfers-dispatch}{Sending a batch onward. The scan box has been
filled by the reader; the note is created when you press the blue button.}{fig:op-dispatch}

The printable note opens in a new tab automatically. Print it and send it with
the trolley.

\begin{warnbox}
If some tags are refused you will see a list with a reason against each one, such
as ``Currently in Finishing, not Sorting Station''. Those garments are not on the
note. Put them aside --- they are in the wrong place and need a supervisor.
\end{warnbox}

\section{Looking back at what your section has moved}

Below the inbox on the \menu{Transfers} screen is \emph{Transfer history for this
section}: every handover your department has been part of, in both directions.

\shot{74-transfers-history}{Transfer history. \texttt{IN} rows are things you
received, \texttt{OUT} rows are things you sent. The \field{Sent} and
\field{Received} columns should match.}{fig:op-history}

\begin{table}[H]
\centering
\small
\begin{tabular}{@{}p{24mm}p{110mm}@{}}
\toprule
\textbf{Column} & \textbf{What it tells you} \\
\midrule
Document & The transfer note number, printed on the paper that travelled with the goods. \\
Route    & Which department sent it and which received it. \\
Batch    & Your own reference for the load, if one was typed in. \\
Sent \newline Received & The two counts. If they differ, the difference is in the next two columns. \\
Missing \newline Extra & How many were on the note but not found, and found but not on the note. \\
Status   & \texttt{RECEIVED} when it is settled, \texttt{VARIANCE} while a difference is still open. \\
\bottomrule
\end{tabular}
\caption{Reading a row of transfer history.}
\end{table}

\begin{tipbox}
Use the \field{Section} box at the top right to look at a different department.
This is the quickest way to answer ``did you ever send us that batch?'' --- the
answer is a document number, a time and a count.
\end{tipbox}

\chapter{The Quality Inspector Portal}
\label{ch:qc}

\portalcover{Quality Inspector Portal}{For staff at the inspection benches}{secQC}

\section{Who this chapter is for}

You inspect finished garments and decide whether each one is good enough to send
to the customer. The system gives you three things the old paper method could
not:

\begin{bullets}
  \item The \textbf{right specification} for the garment in front of you, pulled
        up automatically from its tag.
  \item A picture of the design on which you can mark \textbf{exactly where} the
        fault is.
  \item A record of \textbf{every inspection ever done} on that garment, so you
        know if it has been through before.
\end{bullets}

\section{Your screen}

Open \menu{Quality Control}. There are three tabs across the top.

\begin{table}[H]
\centering
\small
\begin{tabular}{@{}p{32mm}p{106mm}@{}}
\toprule
\textbf{Tab} & \textbf{What it is for} \\
\midrule
Inspect          & Your working screen. Scan a garment and check it. \\
QC queue         & Everything currently sitting in quality control. \\
Defect analysis  & Where faults happen on each design, and which faults are most common. \\
\bottomrule
\end{tabular}
\end{table}

\shot{80-qc-inspect-empty}{The Inspect tab, waiting for you to scan a garment.}{fig:qc-empty}

\section{Inspecting one garment}

\begin{steps}
  \item Put the garment on the bench.
  \item Scan its tag. If your bench reader is set up, the number appears by
        itself; otherwise type it and press Enter.
  \item Press \btn{Look up}.
\end{steps}

The screen now shows everything about this garment.

\shot{82-qc-inspect-garment}{A garment ready for inspection. The design picture is
on the left; the order, customer, size and inspection history on the right.}{fig:qc-garment}

\begin{table}[H]
\centering
\small
\begin{tabular}{@{}p{34mm}p{104mm}@{}}
\toprule
\textbf{What you see} & \textbf{Why it matters} \\
\midrule
Design picture      & The correct design for this garment, front and back. \\
Colour and size     & Check the garment matches. \\
Order and customer  & Different customers accept different tolerances. \\
In QC since         & How long it has been waiting. \\
Previous failures   & If this is not the first time, look harder. \\
\bottomrule
\end{tabular}
\end{table}

\section{Marking a fault on the design}

This is the part that makes the retrofit operator's job possible.

\begin{steps}
  \item Look at the garment and find the fault.
  \item \textbf{Click the same place on the picture.} A small window opens.
  \item Choose the fault from the list, set how serious it is, and add a short
        note if it helps.
  \item Press \btn{Add defect}.
\end{steps}

\shotsmall{83-qc-defect-dialog}{Clicking the design opens this window. Choose what
is wrong and how serious it is.}{fig:qc-dialog}

A numbered marker now sits on the picture exactly where you clicked, and the
fault is listed beside it.

\shot{84-qc-defect-marked}{The fault is marked. Marker 1 on the design is the
same fault as item 1 in the list.}{fig:qc-marked}

\begin{table}[H]
\centering
\small
\begin{tabular}{@{}p{28mm}p{110mm}@{}}
\toprule
\textbf{Severity} & \textbf{Use it when} \\
\midrule
\texttt{CRITICAL} & The garment cannot be sold like this --- an open seam, a hole, a wrong label. \\
\texttt{MAJOR}    & A customer would notice and complain --- a broken stitch, a stain. \\
\texttt{MINOR}    & Cosmetic; a customer probably would not notice --- slight puckering. \\
\bottomrule
\end{tabular}
\caption{Choosing the right severity.}
\end{table}

\begin{tipbox}
\textbf{Use both views.} Press \btn{Front} or \btn{Back} above the picture to mark
faults on either side. If a fault has no particular place --- a measurement that
is out of tolerance, for instance --- use \btn{+ Add defect without a position}.
\end{tipbox}

\begin{tipbox}
Clicked in the wrong spot? Just click the marker itself and it disappears.
\end{tipbox}

\section{Passing or failing}

Two large buttons sit at the top right of the garment panel.

\begin{table}[H]
\centering
\small
\begin{tabular}{@{}p{38mm}p{100mm}@{}}
\toprule
\textbf{Button} & \textbf{What happens next} \\
\midrule
\btn{Pass} &
  The garment is cleared and can go to dispatch. \\
\btn{Fail --- send to retrofit} &
  A repair job is opened automatically, carrying every marker you placed.
  The garment then needs sending to the retrofit bench on a transfer note. \\
\bottomrule
\end{tabular}
\end{table}

\begin{stopbox}
\textbf{You cannot fail a garment without giving a reason.} If you press
\btn{Fail} with no faults marked, the system stops you. This is deliberate: the
retrofit operator has to know what to fix.
\end{stopbox}

After you decide, the screen offers \btn{Inspect the next garment} so you can
carry straight on.

\section{Passing a whole clean batch}

When you have checked a pile by hand and every piece is good, you do not have to
scan them one at a time.

\begin{steps}
  \item Scroll to \emph{Pass a whole batch} at the bottom of the Inspect tab.
  \item Bulk-read the whole pile into the scan box.
  \item Press \btn{Pass all scanned garments}.
\end{steps}

\begin{warnbox}
Every pass is still recorded individually against \emph{your} name. Only use the
batch button when you really have inspected every garment in the pile.
\end{warnbox}

Anything that could not be passed --- because it is not in quality control, for
example --- is listed afterwards so nothing is silently skipped.

\section{The QC queue}

The \menu{QC queue} tab shows everything in the department, oldest first.

\shot{81-qc-queue}{The QC queue. Click any row to inspect that garment.}{fig:qc-queue}

The four boxes at the top tell you at a glance:

\begin{bullets}
  \item how many are \textbf{awaiting inspection},
  \item how many have \textbf{passed} and are waiting to go to dispatch,
  \item how many have \textbf{failed} and are waiting to go to retrofit,
  \item how many \textbf{batches are inbound} from finishing.
\end{bullets}

Clicking any row takes you straight to the inspection screen for that garment.

\section{Defect analysis}

The \menu{Defect analysis} tab answers the question ``where do our faults keep
happening?''

\shot{85-qc-analysis}{Defect analysis. Every marker ever placed on this design is
shown on one picture, so patterns stand out.}{fig:qc-analysis}

\begin{bullets}
  \item The picture on the left shows \textbf{every fault ever recorded} on that
        design. A cluster of markers in one place usually means a machine or an
        operation needs attention, not a careless worker.
  \item The table on the right ranks faults by how often they occur.
  \item The bottom table shows inspector activity.
\end{bullets}

\begin{plainwords}
A high fail rate for one inspector does not automatically mean they are wrong ---
they may simply be checking the hardest styles. Use this table to start a
conversation, never to draw a conclusion on its own.
\end{plainwords}

\section{Receiving and sending batches}

Quality control receives from finishing and sends to either retrofit or dispatch.
Both work exactly as described in Section~\ref{sec:receive}
and Section~\ref{sec:dispatch}.

\begin{tipbox}
When sending passed garments to dispatch, the system checks each one really has
passed. Anything that has not is refused with a clear reason, so a failed garment
can never slip into a customer shipment.
\end{tipbox}


\chapter{Glossary}
\label{ch:glossary}

Every term used in the software and in this manual, in alphabetical order and in
plain language.

\begin{longtable}{@{}p{40mm}p{98mm}@{}}
\toprule
\textbf{Term} & \textbf{Meaning} \\
\midrule
\endfirsthead
\toprule
\textbf{Term} & \textbf{Meaning} \\
\midrule
\endhead
\bottomrule
\endfoot

Article &
  One single garment. Every article has its own tag and its own serial number. \\

Audit trail &
  The permanent record of who changed what and when. It cannot be edited. \\

Batch &
  A group of garments moved together, normally all the same design and size. \\

Batch reference &
  A name you give a batch, such as \typed{WASH-LOT-07}, printed on the transfer
  note so the receiving department can find it. \\

Bundle &
  A tied stack of cut pieces handed from cutting to a sewing line. \\

Commissioning &
  Attaching a tag to a garment and registering it. After this the garment exists
  as far as the system is concerned. \\

Cut order &
  An instruction to cut a particular design and colour in a particular quantity. \\

Defect map &
  The picture of the design that quality control clicks on to show exactly where
  a fault is. \\

Dispatch (verb) &
  To send goods from one department to the next, creating a transfer note. \\

Dispatch (department) &
  The final section, where the factory tag is removed and the customer's tag is
  applied. \\

EPC &
  The unique number stored inside an RFID tag. Short for ``Electronic Product
  Code''. You will see it as a long string of letters and digits. \\

First-pass yield &
  The percentage of garments that passed quality control the very first time,
  without needing any correction. \\

GRN &
  ``Goods Receipt Note'' --- the record created when fabric arrives from a
  supplier. \\

Handheld / wedge mode &
  A hand scanner set up to behave like a keyboard, so the tags it reads simply
  type themselves into whatever box the cursor is in. \\

In transit &
  A garment that has been dispatched by one department but not yet received by
  the next. It belongs to neither at that moment. \\

Master data &
  The reference lists --- designs, colours, sizes, customers, fabric types,
  defect codes --- that everything else refers to. \\

On hold &
  A garment taken out of the normal flow, usually after a variance was closed. It
  needs a supervisor to release it. \\

Portal (reader) &
  A doorway-shaped reader that senses tags as goods pass through it. \\

Retrofitting &
  The department that corrects faults found by quality control. \\

Roll &
  A roll of denim fabric. May carry its own tag. \\

Scrap &
  A garment written off as unsaleable. It leaves work in process permanently. \\

Serial number &
  A readable name for a garment, such as \typed{ART-260816-000123}. Easier to say
  aloud than the tag number. \\

Severity &
  How serious a fault is: \texttt{CRITICAL}, \texttt{MAJOR} or \texttt{MINOR}. \\

Shift &
  The working period. This factory runs two eight-hour shifts, A and B. \\

Sorting session &
  A working period at a sorting table, during which a pile is read and split into
  groups. \\

Style &
  A design, such as ``Slim Fit 5-Pocket Jean''. Carries the pictures used by
  quality control. \\

Tag &
  The small electronic label attached to a garment or roll. Contains no battery. \\

Tag swap &
  Removing the factory tracking tag at dispatch and applying the customer's own
  tag in its place. \\

Tally &
  Comparing what was received against what was sent. \\

Transfer note &
  The document created when goods move between departments, listing exactly what
  was sent. \\

Tunnel (reader) &
  A conveyor-based reader. The best way to read a dense pile reliably. \\

Unbound tag &
  A tag that has been released from a garment. It no longer identifies anything
  and can be reused on a new garment. \\

Variance &
  The count on arrival did not match the count that was sent. \\

WIP &
  ``Work in progress'' --- everything on the factory floor that has not yet been
  shipped. \\

\end{longtable}

\chapter{When Something Goes Wrong}
\label{ch:trouble}

The system refuses things on purpose, to protect the count. Nearly every message
below is the software doing its job. Find the message you saw, and follow the
advice.

\section{Signing in}

\begin{longtable}{@{}p{52mm}p{86mm}@{}}
\toprule
\textbf{Message} & \textbf{What to do} \\
\midrule
\endfirsthead
\toprule
\textbf{Message} & \textbf{What to do} \\
\midrule
\endhead
\bottomrule
\endfoot

Invalid username or password &
  Check Caps~Lock. Check you are using your own account. After a few tries, ask
  your supervisor to reset it. \\

Your session has expired &
  Normal after twelve hours, or if you were signed in elsewhere. Sign in again.
  Nothing you saved is lost. \\

Cannot reach the server &
  The network or the server is down. Tell your supervisor. Do not keep working on
  paper without telling anyone --- the count will drift. \\

\end{longtable}

\section{Scanning and receiving}

\begin{longtable}{@{}p{52mm}p{86mm}@{}}
\toprule
\textbf{Message} & \textbf{What it means and what to do} \\
\midrule
\endfirsthead
\toprule
\textbf{Message} & \textbf{What it means and what to do} \\
\midrule
\endhead
\bottomrule
\endfoot

Tag is not registered to any article &
  This tag was never commissioned, or it was released at dispatch and is now
  blank. Set the garment aside for a supervisor. \\

Currently in \emph{[department]}, not \emph{[your department]} &
  The system believes this garment is somewhere else. Usually a trolley moved
  without being scanned. Set it aside and tell your supervisor, who can correct
  the record. \\

Already dispatched on \emph{[document]} &
  Somebody has already sent this garment. Check whether two people scanned the
  same trolley. \\

Article is on hold &
  It was written off in an earlier variance. A supervisor must release it. \\

Count does not match &
  Not an error. Search the trolley and scan again --- anything found is added and
  the variance clears. \\

Nothing scanned &
  The scan box is empty. Click inside it before pulling the trigger. \\

\end{longtable}

\section{Attaching tags at stitching}

\begin{longtable}{@{}p{52mm}p{86mm}@{}}
\toprule
\textbf{Message} & \textbf{What it means and what to do} \\
\midrule
\endfirsthead
\toprule
\textbf{Message} & \textbf{What it means and what to do} \\
\midrule
\endhead
\bottomrule
\endfoot

\emph{n} tag(s) are already registered to other articles &
  Those tags are in use on other garments. The list is shown. Take them off, use
  fresh tags, and hand the used ones to your supervisor. \\

Bundle has \emph{n} garment(s) left to tag but \emph{m} tags were scanned &
  You have more tags than pieces. Either you scanned a tag twice from a nearby
  tray, or the bundle count is wrong. Re-count the bundle. \\

Bundle has not been issued to stitching yet &
  Cutting has not handed it over. Ask them to press
  \btn{Issue all to stitching}. \\

\end{longtable}

\section{Quality control}

\begin{longtable}{@{}p{52mm}p{86mm}@{}}
\toprule
\textbf{Message} & \textbf{What it means and what to do} \\
\midrule
\endfirsthead
\toprule
\textbf{Message} & \textbf{What it means and what to do} \\
\midrule
\endhead
\bottomrule
\endfoot

A QC failure must record at least one defect &
  You pressed Fail without marking anything. Click the design where the fault is
  and choose a defect first. \\

\emph{[serial]} is in \emph{[department]}, not QC &
  The garment has not been received into quality control. Receive the batch on
  the Transfers screen first. \\

The picture does not match the garment &
  The design image in Master Data is out of date. Do not mark faults on the wrong
  picture --- tell your administrator. \\

\end{longtable}

\section{Retrofitting}

\begin{longtable}{@{}p{52mm}p{86mm}@{}}
\toprule
\textbf{Message} & \textbf{What it means and what to do} \\
\midrule
\endfirsthead
\toprule
\textbf{Message} & \textbf{What it means and what to do} \\
\midrule
\endhead
\bottomrule
\endfoot

Describe the correction &
  The \field{Correction carried out} box is empty. Write what you actually did. \\

No open retrofit job for this article &
  This garment has not been failed by quality control, or the job is already
  closed. Check the serial number. \\

\emph{[serial]} is in \emph{[department]}, not Retrofitting &
  Receive the batch into retrofit before completing the job. \\

\end{longtable}

\section{Dispatch}

\begin{longtable}{@{}p{52mm}p{86mm}@{}}
\toprule
\textbf{Message} & \textbf{What it means and what to do} \\
\midrule
\endfirsthead
\toprule
\textbf{Message} & \textbf{What it means and what to do} \\
\midrule
\endhead
\bottomrule
\endfoot

\emph{[serial]} has not passed QC &
  Only inspected and passed garments can be re-tagged. Send it back to quality
  control. \\

Customer tag \ldots\ is already in use &
  That customer tag has been applied to another garment. Use a fresh one. \\

\emph{[serial]} already carries customer tag \ldots &
  This garment has already been re-tagged. Do not do it twice. \\

The customer tag must differ from the tracking tag &
  You scanned the same tag into both boxes. Scan the new tag into the second
  box. \\

\end{longtable}

\section{Reading tags reliably}

If a reader keeps missing tags, the cause is almost always physical rather than
software.

\begin{table}[H]
\centering
\small
\begin{tabular}{@{}p{46mm}p{92mm}@{}}
\toprule
\textbf{Symptom} & \textbf{Usual cause} \\
\midrule
A few tags missed in a dense pile &
  Tags are shielded by other garments. Move the pile through the reader rather
  than holding it still, or read it in two halves. \\

Nothing reads at all &
  The reader has lost power or network. Check \field{Last seen} on the readers
  screen. \\

Tags read from the next station &
  The reader's power is set too high, or antennas are pointing the wrong way.
  Report it --- stray reads cause wrong-section errors. \\

One particular garment never reads &
  Its tag is damaged. A supervisor can fit a replacement without losing the
  garment's history. \\

\bottomrule
\end{tabular}
\end{table}

\section{Who to ask}

\begin{table}[H]
\centering
\small
\begin{tabular}{@{}p{46mm}p{92mm}@{}}
\toprule
\textbf{Problem} & \textbf{Ask} \\
\midrule
A garment is in the wrong section        & Your supervisor \\
Missing pieces after two scans           & Your supervisor \\
A tag will not read                      & Your supervisor \\
You cannot sign in                       & Your supervisor, then the administrator \\
A reader is offline                      & The administrator \\
A design picture is wrong                & The administrator \\
You need a new kind of report            & Your supervisor or the plant manager \\
\bottomrule
\end{tabular}
\end{table}


\end{document}
